\section{Cohomology of intersection-closed families (UC9--UC12)}
\label{sec:cohomology}

The cohomological arc works in the intersection-closed dialect
(Section~\ref{ssec:flip}). Its aim is to attach to a hypothetical
counterexample a cohomology class --- the \emph{obstruction class} ---
that two independent arguments force to be simultaneously nonzero and
zero. This section builds the cohomology theory the obstruction lives
in (UC9--UC10) and proves the one genuinely unconditional theorem of
the arc, the trace-injectivity / null-homotopy result (UC12).
Section~\ref{sec:obstruction} then assembles the obstruction class
itself.

\subsection{The U1 wall: why a direct functorial proof fails (UC9)}
\label{ssec:u1-wall}

A natural first attempt embeds intersection-closed families into a
combinatorial object whose cohomology is already known. UC9 carries
this out. The sibling $1/3$--$2/3$ program supplies, as an external
input, a computation of the rational cohomology of the order complex
$\Delta(\mathrm{PPF}_n)$ of proper partial orders on $[n]$:

\begin{remark}[External input]
\label{rem:external-input}
The $1/3$--$2/3$ program asserts $\widetilde
H^{k}(\Delta(\mathrm{PPF}_n);\Q)\cong\sgn_{\Sym_n}$ for $k=n-2$ and
$0$ otherwise --- the rational cohomology of an $(n-2)$-sphere with the
sign action. This is consumed here as a black box. A referee assessing
the cohomological arc must check it against its stated source; the
present paper does not re-derive it.
\end{remark}

UC9 constructs a \emph{cone embedding} $\iota^{\mathrm{cone}}_n\colon
2^{[n]}\setminus\{\emptyset\}\hookrightarrow\mathrm{PPF}_{n+1}$, $S\mapsto
\omega_S$, where $\omega_S$ is the partial order on $[n+1]$ placing
$n+1$ below exactly the elements of $S$ with $[n]$ an antichain. The
embedding is order-, meet- and join-preserving, $\Sym_n$-equivariant,
and injective. It is, in a precise sense, the unique such embedding.

The attempt then fails, and the way it fails is instructive. The
embedded image $\Delta(2^{[n]}\setminus\{\emptyset\})$ is
\emph{contractible} (a cone with apex $[n]$). Hence the induced
pullback on cohomology is the zero map in every degree, and the
relative cohomology $\widetilde H^*(\Delta_{n+1},\Delta(\mathrm{im}))$
is \emph{family-independent}: it sees nothing of the particular
intersection-closed family. A family-independent invariant cannot
distinguish two families, hence cannot detect a counterexample. UC9
names this the \emph{U1 wall} and records its verdict honestly as
\emph{partial}: the cohomology pulls back cleanly, but Frankl-extraction
fails because the pullback is trivial.

\begin{remark}[The lesson of the U1 wall]
\label{rem:u1-lesson}
The U1 wall is the reason the program does not pursue a direct
functorial proof. Any argument of the shape ``embed families into a
fixed space, read off a cohomological invariant'' is killed by
contractibility of the image. The cohomological arc therefore builds a
\emph{native, family-internal} cohomology theory instead
(Section~\ref{ssec:cn-category}), in which the family is part of the
indexing data rather than a sub-object of a fixed ambient space.
\end{remark}

\subsection{The category of intersection-closed families (UC10)}
\label{ssec:cn-category}

\begin{definition}
\label{def:cn-category}
The category $\Cn$ has as objects the pointed intersection-closed
families: pairs $(S,\F)$ with $S\subseteq[n]$, $\F\subseteq2^S$
intersection-closed, $\bigcup\F=S$, and $S\in\F$. A morphism
$(S,\F)\to(T,\G)$ exists when $T\subseteq S$ and $\G=\{A\cap T:A\in\F\}$
is the \emph{trace} of $\F$ to $T$; trace preserves
intersection-closure, so this is well-defined. Same-ground-set
inclusions are deliberately \emph{excluded} --- only traces are
morphisms.
\end{definition}

The symmetry group is the hyperoctahedral (type-$B_n$ Coxeter) group
\[
  \Gamma_n=(\Z/2)^n\rtimes\Sym_n,\qquad|\Gamma_n|=2^n\,n!,
\]
with $(\Z/2)^n$ generated by the coordinate toggles $\sigma_x$
(Section~\ref{ssec:walsh}) and $\Sym_n$ relabelling coordinates.
Centring $(\Z/2)^n$ --- not $\Sym_n$ --- is the decisive design choice:
it makes the Walsh decomposition, not the Specht/FI machinery of the
$1/3$--$2/3$ side, the organising tool.

\begin{definition}
\label{def:cubical-defect-complex}
For a pointed family $(S,\F)$ the \emph{cubical Walsh-defect complex}
$X(\F)$ is the cubical complex whose $k$-cells are the $k$-dimensional
subcubes of $Q_S$ all of whose $2^k$ vertices lie in $\F$. The program's
cohomology object is the homotopy colimit
\[
  \Xn\;=\;\hocolim_{(\Cn)^{\mathrm{op}}}X(\F),
\]
a $\Gamma_n$-equivariant cochain complex over $\Q$, presented
concretely by the Bousfield--Kan double complex~\cite{bousfield-kan}
\[
  \BK^{p,q}(\Xn)\;=\;\prod_{\F_0\to\F_1\to\cdots\to\F_p\ \text{in }\Cn}
  C^q\bigl(X(\F_p);\Q\bigr),
\]
with horizontal (bar / cosimplicial) differential $\delta_{\BK}$ and
vertical (cubical cochain) differential $d$; $C^*(\Xn;\Q)$ is its total
complex.
\end{definition}

\begin{theorem}[Walsh decomposition, UC10]
\label{thm:walsh-decomposition}
The differentials of $C^*(\Xn;\Q)$ are $(\Z/2)^n$-equivariant. Hence,
by Maschke's theorem, the cochain complex splits canonically as a
direct sum of sub-cochain-complexes indexed by Walsh characters,
\[
  C^*(\Xn;\Q)\;=\;\bigoplus_{S\subseteq[n]}\walsh S\otimes
  \isotype S{*},
\]
and $\Sym_n$ permutes the summands by $\isotype S*\mapsto
\isotype{\pi(S)}*$. We call $\isotype Sk$ the
\emph{$\walsh S$-isotype multiplicity complex in degree $k$}, and
$|S|$ its Walsh level.
\end{theorem}

This is solid: it is just Maschke's theorem applied degreewise to a
finite-group action that commutes with the differential. It is the
foundation on which both $Y_1$ and $Y_2$ are stated, and it is not in
dispute.

The cohomology object is linked back to the bias by the
\emph{stuck-set correspondence}: for a single-coordinate trace
$(S,\F)\to(S\setminus\{x\},\G)$ each $A\in\G$ has one or two preimages,
and the count of two-preimage versus one-preimage members reproduces
$\bias_x(\F)=|\F^{\stuck}_x|-|\F^{\stuck}_{\bar x}|$. This is the
bridge from cohomology to Frankl; it is used in
Section~\ref{sec:obstruction}.

\subsection{The sphere-concentration target UC10.1}
\label{ssec:uc10-1}

The program's structural goal is that $\Xn$ is, cohomologically, a
sphere with the top Walsh / sign action --- the cube-side analogue of
the input in Remark~\ref{rem:external-input}.

\begin{conjecture}[UC10.1, sphere concentration --- the program target]
\label{conj:uc10-1}
For $n\ge3$,
\[
  \widetilde H^{k}(\Xn;\Q)\;\cong\;
  \begin{cases}
    \walsh{[n]}\boxtimes\sgn_{\Sym_n}, & k=n-1,\\[2pt]
    0, & 1\le k\le n-2,
  \end{cases}
\]
as $\Gamma_n$-modules. Equivalently: the top isotype
$\isotype{[n]}{n-1}$ is one-dimensional and every other isotype
$\isotype Sk$ with $k\ge1$ vanishes.
\end{conjecture}

\begin{status}
Conjecture~\ref{conj:uc10-1} is \emph{stated, not proved}. The original
UC10 note records it as the target with verdict \AMBER{partial}. It
decomposes into:
\begin{itemize}[leftmargin=1.4em]
\item \textbf{the trace-injectivity input} --- proved unconditionally,
  Theorem~\ref{thm:trace-injectivity} below (UC12);
\item \textbf{the top-Walsh concentration} ($\isotype{[n]}k=0$ for
  $1\le k\le n-2$) --- argued in UC13~\S5 / UC14~$R3$ by a
  cofiber long-exact-sequence and induction; the \emph{method} is
  sound, with soft spots noted (Section~\ref{ssec:y1});
\item \textbf{the lower-Walsh vanishing} $Y_1$ ($\isotype Sk=0$ for
  $S\subsetneq[n]$, $k\ge1$) --- \emph{open}; the previously
  circulated proof is invalid (Correction~\ref{corr:trace-exact}).
\end{itemize}
Thus Conjecture~\ref{conj:uc10-1} is not available as a theorem; the
honest status of $Y_1$ is Section~\ref{ssec:y1}.
\end{status}

\subsection{The doubling functor and trace-injectivity (UC12)}
\label{ssec:uc12}

This subsection states and proves the one genuinely unconditional
theorem of the cohomological arc. It is the UC12 result, and it does
not depend on any disputed step.

\begin{definition}
\label{def:doubling}
The \emph{doubling functor} $\mathrm{db}\colon\Cn\to\mathcal
C^{\cap}_{n+1}$ sends $(S,\F)$ to
$(S\sqcup\{n+1\},\ \F\cup(\F+\{n+1\}))$ where $\F+\{n+1\}=\{A\cup\{n+1\}
:A\in\F\}$. It is fully faithful, and geometrically
$X(\mathrm{db}\,\F)\cong X(\F)\times I$ --- the \emph{matched cylinder}
(prism): every $k$-cell of $X(\F)$ acquires a \emph{bridge}
$(k+1)$-cell, its prism in the new coordinate. The canonical inclusion
$\iota^{\cap}_n\colon\Xn\hookrightarrow X^{\cap}_{n+1}$ and the trace
projection $\rho_n$ make $\mathrm{db}$ a section of $\rho_n$.
\end{definition}

\begin{construction}[The cubical-bridge null-homotopy]
\label{constr:bridge}
On the top isotype $\isotype{[n+1]}*$ define $h\colon
C^{k}(X^{\cap}_{n+1})\to C^{k-1}(\Xn)$ by $(h\omega)(D)=\omega
(\widetilde D)$, where $\widetilde D=D\times\{0\!\to\!1\}$ is the
bridge cell of $D$.
\end{construction}

\begin{theorem}[Trace-injectivity / chain-homotopy identity, UC12]
\label{thm:trace-injectivity}
For every cochain $\omega\in\isotype{[n+1]}{k}(X^{\cap}_{n+1})$,
\[
  (\iota^{\cap}_n)^{*}\omega\;=\;\frac{(-1)^{k}}{2}\,
  \bigl(d\,h\,\omega-h\,d\,\omega\bigr).
\]
Consequently $(\iota^{\cap}_n)^{*}=0$ on $\isotype{[n+1]}{}$-cohomology,
and the connecting map $\delta^{\cap}_n$ of the cofiber long exact
sequence is \emph{injective} on the top-isotype piece. The statement
holds on the full bi-isotype $\walsh{[n+1]}\boxtimes\sgn_{\Sym_{n+1}}$.
\end{theorem}

\begin{proof}[Proof sketch]
The prism boundary formula gives $\partial\widetilde C=\sum_F[F:C]
\widetilde F+(-1)^{k}(C\times\{1\}-C\times\{0\})$ for a $k$-cell $C$.
Pair with $\omega$ and use \emph{top-Walsh antisymmetry}: on the
$\walsh{[n+1]}$-isotype the toggle $\sigma_{n+1}$ acts by $-1$, so
$\omega(C\times\{1\})=(\sigma_{n+1}^{*}\omega)(C\times\{0\})=
-\omega(C\times\{0\})$. The two boundary-face terms combine to
$2\omega(C\times\{0\})=2\,((\iota^{\cap}_n)^{*}\omega)(C)$ up to sign,
and the bar-string / hocolim coherence of the prism is exactly the
content of $\mathrm{db}$ being a fully faithful section. The
chain-homotopy identity follows; for a cocycle $\omega$ it shows the
trace is a coboundary, hence $(\iota^{\cap}_n)^{*}=0$ on cohomology;
injectivity of $\delta^{\cap}_n$ is then immediate from the cofiber
sequence.
\end{proof}

\begin{status}
Theorem~\ref{thm:trace-injectivity} is \emph{genuine and
unconditional}. It is the UC12 result; the matched-cylinder prism
structure is real, and the antisymmetry step is forced by the
definition of the top isotype. Its machine-checked formalisation is
clean (no axiom, no \texttt{sorry}); see Section~\ref{ssec:lean-state}.
This is the cohomological arc's solid ground. Everything in
Section~\ref{sec:obstruction} and the audit of
Section~\ref{sec:proof-state} should be read against it: the
trace-injectivity theorem is \emph{not} where the program's difficulty
lies.
\end{status}

\begin{remark}
\label{rem:bridge-vs-y1}
The bridge of Theorem~\ref{thm:trace-injectivity} gives exactly
$(\iota^{\cap}_n)^{*}=0$ on cohomology --- the trace map is
null-homotopic. It does \emph{not} give vanishing of the source
cohomology. This distinction --- ``the trace is exact'' versus ``the
class is zero'' --- is benign here, because
Theorem~\ref{thm:trace-injectivity} only ever claims the former. It
becomes the central error of the $Y_1$ proof
(Correction~\ref{corr:trace-exact}), which silently upgrades the former
to the latter. The honest use of a null-homotopic trace map is to feed
it into a long exact sequence; Section~\ref{ssec:y1} does exactly that.
\end{remark}
